% TEMPLATE-MILC-v3.tex - plantilla maestra para todas las guias MILC v3
% Ver scripts/generadores/build-guia-milc.py para la lista completa de
% claves esperadas. El script valida que no queden placeholders sin
% reemplazar antes de escribir el archivo final.

\documentclass[11pt,letterpaper]{article}
\usepackage[margin=1.75cm]{geometry}
\usepackage[table]{xcolor}
\usepackage{fontspec}
\usepackage[spanish]{babel}
\usepackage{tikz}
\usepackage[most]{tcolorbox}
\usepackage{tabularx}
\usepackage{array}
\usepackage{fancyhdr}
\usepackage{titlesec}
\usepackage{ragged2e}
\usepackage{enumitem}
\usepackage{microtype}

\setmainfont{Helvetica Neue}
\setsansfont{Helvetica Neue}
\setlength{\parindent}{0pt}
\setlength{\parskip}{6pt}
\hyphenpenalty=200
\exhyphenpenalty=200
\pretolerance=400
\tolerance=2000
\setlength{\emergencystretch}{1em}
\renewcommand{\arraystretch}{1.25}
\setlist{leftmargin=5mm,itemsep=1mm,topsep=1mm}
\newcolumntype{Y}{>{\RaggedRight\arraybackslash}X}

\definecolor{milcMagenta}{HTML}{E6007E}
\definecolor{milcVino}{HTML}{5A0038}
\definecolor{milcTurquesa}{HTML}{0093A5}
\definecolor{milcVerde}{HTML}{58A923}
\definecolor{milcMostaza}{HTML}{E5B400}
\definecolor{milcOcre}{HTML}{B86600}
\definecolor{milcNegro}{HTML}{241816}
\definecolor{milcGris}{HTML}{F7F7F4}
\definecolor{milcLinea}{HTML}{E8E2DD}
\definecolor{milcGradePrimary}{HTML}{0066FF}
\definecolor{milcGradeDark}{HTML}{003D99}
\definecolor{milcGradeSoft}{HTML}{D6E8FF}
\definecolor{milcGradeAccent}{HTML}{CFE7FF}
\definecolor{milcGradeText}{HTML}{FFFFFF}
\color{milcNegro}

\pagestyle{fancy}
\fancyhf{}
\lhead{\small Tecnología e Informática · Grado 11}
\rhead{\small Guía 3 · MILC v3}
\cfoot{\small\thepage}
\renewcommand{\headrulewidth}{0pt}

\newtcolorbox{titlebox}[1]{
  enhanced,colback=#1,colframe=#1,arc=7mm,boxrule=0pt,
  left=7mm,right=7mm,top=3mm,bottom=3mm,
  before skip=8pt,after skip=8pt,
  fontupper=\sffamily\bfseries\Large\color{white}
}

\newtcolorbox{softbox}[3]{
  enhanced,colback=#2,colframe=#1,borderline west={4pt}{0pt}{#1},
  arc=2mm,boxrule=.4pt,left=4mm,right=4mm,top=3mm,bottom=3mm,
  before skip=6pt,after skip=8pt,title={#3},coltitle=white,
  colbacktitle=#1,fonttitle=\sffamily\bfseries,fontupper=\RaggedRight,
  attach boxed title to top left={xshift=3mm,yshift=-2mm},
  boxed title style={arc=2mm, boxrule=0pt}
}

\newtcolorbox{drawbox}[2]{
  enhanced,colback=white,colframe=#1,arc=4mm,boxrule=1pt,
  left=5mm,right=5mm,top=4mm,bottom=4mm,before skip=8pt,after skip=8pt,
  title={#2},coltitle=white,colbacktitle=#1,fonttitle=\sffamily\bfseries,
  fontupper=\RaggedRight,
  attach boxed title to top left={xshift=4mm,yshift=-2mm},
  boxed title style={arc=3mm, boxrule=0pt}
}

\newtcolorbox{infoband}[1]{
  enhanced,colback=#1!8,colframe=#1!50,arc=2mm,boxrule=.5pt,
  left=4mm,right=4mm,top=2mm,bottom=2mm,before skip=4pt,after skip=4pt,
  fontupper=\small\RaggedRight
}

% Puente narrativo entre secciones (contrato editorial Conéctate).
% Da continuidad: "Acabas de X. Ahora vas a Y porque Z."
% Visualmente: banda izquierda en color del grado, texto en itálica pequeña.
\newtcolorbox{puentebox}{
  enhanced,colback=milcGradeSoft,colframe=milcGradeDark,
  borderline west={3pt}{0pt}{milcGradeDark},
  arc=0mm,boxrule=0pt,left=5mm,right=4mm,top=2.5mm,bottom=2.5mm,
  before skip=4pt,after skip=8pt,
  fontupper=\itshape\small\color{milcNegro}\RaggedRight
}
\newcommand{\puente}[1]{%
  \begin{puentebox}%
    \textcolor{milcGradeDark}{\textbf{\textrightarrow}}\hspace{2mm}#1%
  \end{puentebox}%
}

\newcommand{\linead}[1]{\noindent\color{milcLinea}\rule{#1}{0.5pt}\color{milcNegro}}
\newcommand{\respuesta}[1]{\par\vspace{2mm}\foreach \n in {1,...,#1}{\noindent\color{milcLinea}\rule{\linewidth}{0.5pt}\par\vspace{5mm}}\color{milcNegro}}
\newcommand{\checkbox}{\textcolor{milcGradeDark}{\fbox{\phantom{x}}}\hspace{5pt}}

\newcommand{\phasecard}[4]{
\begin{tcolorbox}[enhanced,colback=#3,colframe=#2,arc=3mm,boxrule=.5pt,height=3.05cm,valign=center,halign=center,left=1.5mm,right=1.5mm,top=2mm,bottom=2mm]
{\sffamily\bfseries\fontsize{22}{24}\selectfont\textcolor{#2}{#1}}\par
{\sffamily\bfseries\small #4}
\end{tcolorbox}
}

\begin{document}
\thispagestyle{empty}
\begin{tikzpicture}[remember picture,overlay]
  \fill[white] (current page.south west) rectangle (current page.north east);
  \path[fill=milcGradePrimary,rounded corners=16mm]
    ([xshift=.55cm,yshift=-.55cm]current page.north west) rectangle
    ([xshift=-.55cm,yshift=3.65cm]current page.south east);

  \node[anchor=north west,text=milcGradeText] at ([xshift=2.0cm,yshift=-1.85cm]current page.north west)
    {\sffamily\bfseries\fontsize{14}{16}\selectfont GRADO};
  \node[anchor=north west,text=milcGradeText,opacity=.82] at ([xshift=2.0cm,yshift=-2.75cm]current page.north west)
    {\sffamily\bfseries\fontsize{9}{11}\selectfont SERIE GUÍAS · TECNOLOGÍA E INFORMÁTICA};

  \begin{scope}[shift={([xshift=-3.05cm,yshift=-2.55cm]current page.north east)}]
    \fill[white,opacity=.80] (-.62,-.52) rectangle (.62,.52);
    \draw[milcGradeText,line width=1pt] (-.62,-.52) rectangle (.62,.52);
    \fill[milcGradeDark] (-.42,-.38) rectangle (-.22,.06);
    \fill[milcGradeAccent] (-.10,-.38) rectangle (.10,.24);
    \fill[milcGradePrimary!60!black] (.22,-.38) rectangle (.42,.38);
  \end{scope}

  \node[anchor=west,text=milcGradeText] at ([xshift=1.9cm,yshift=-12.85cm]current page.north west)
    {\sffamily\bfseries\fontsize{124}{124}\selectfont 11};
  \draw[milcGradeText,line width=1.7pt]
    ([xshift=4.52cm,yshift=-11.68cm]current page.north west) circle (.13cm);

  \node[anchor=west,text=milcGradeText,text width=8.7cm,align=left] at ([xshift=10.35cm,yshift=-11.6cm]current page.north west)
    {
     {\sffamily\bfseries\fontsize{19}{22}\selectfont Sitio web personal ---\\fundamentos de\\la casa digital}\\[4mm]
     {\sffamily\bfseries\fontsize{23}{26}\selectfont Guía 3-11-TIC}\\[2mm]
     {\sffamily\fontsize{13}{16}\selectfont Período 1 · Presencia y marca digital}\\[5mm]
     {\sffamily\bfseries\fontsize{11}{13}\selectfont Institución Educativa Sor María Juliana}\\[1mm]
     {\sffamily\fontsize{10}{12}\selectfont Autor: Dr. Álvaro Cárdenas Orozco}
    };

  \path[fill=milcGradeSoft,rounded corners=7mm]
    ([xshift=1.35cm,yshift=.85cm]current page.south west) rectangle
    ([xshift=-1.35cm,yshift=3.15cm]current page.south east);
  \draw[milcGradeDark,line width=.65pt,rounded corners=7mm]
    ([xshift=1.35cm,yshift=.85cm]current page.south west) rectangle
    ([xshift=-1.35cm,yshift=3.15cm]current page.south east);
  \node[anchor=center,text=milcNegro,text width=17.0cm,align=center] at ([yshift=2.0cm]current page.south)
    {\sffamily\fontsize{10}{12}\selectfont Metodología MILC · Apertura ancestral · Pensamiento computacional · Triángulo Dussel-Estoicismo-Floridi\\
    \textcolor{milcGradeDark}{\bfseries Producto:} Wireframe a mano + archivo index.html con estructura HTML semántica};
\end{tikzpicture}
\mbox{}
\newpage

\section*{Apertura: Sitio web personal --- fundamentos de la casa digital}

\begin{softbox}{milcGradeDark}{milcGradeSoft}{Saber ancestral}
Los \textbf{caminos prehispánicos} (la red del oro Quimbaya en el Valle, los caminos incas, los senderos de los chasquis) no eran trazos al azar: cada uno tenía un propósito y conectaba un lugar específico con otro. El que caminaba sabía adónde iba porque la ruta estaba pensada. Esa misma intención organiza un buen sitio web.
\end{softbox}

\begin{softbox}{milcTurquesa}{milcGris}{Saber contemporáneo}
Una página web bien construida usa \textbf{HTML semántico}: \texttt{header}, \texttt{nav}, \texttt{main}, \texttt{article}, \texttt{aside}, \texttt{footer}. Cada etiqueta no solo dibuja: \emph{le dice al navegador} qué función cumple ese bloque. Sin semántica, los lectores de pantalla, los buscadores y las personas se pierden.
\end{softbox}

\begin{softbox}{milcMostaza}{milcGris}{Pregunta-puente}
\textit{¿Qué hacía que los caminos antiguos llevaran a la gente al lugar correcto sin GPS? ¿Qué señales usaban y por qué funcionaban?}
\end{softbox}

\begin{softbox}{milcVerde}{milcGris}{Saber y saber hacer}
Al terminar podrás: (1) \textbf{identificar} la estructura semántica de sitios reales; (2) \textbf{explicar} qué función cumple cada etiqueta semántica; (3) \textbf{crear} el wireframe y el HTML de tu página personal; (4) \textbf{evaluar} la estructura del sitio de un compañero.
\end{softbox}

\small
\begin{tabularx}{\textwidth}{>{\bfseries}p{3.0cm}Y}
\rowcolor{milcGradePrimary!12}Autor & Dr. Álvaro Cárdenas Orozco\\
DBA & Comunicación profesional en entornos digitales (MEN, Lineamientos T\&I)\\
Referentes & Floridi (infoética) · Dussel (estética de la liberación) · Estoicismo\\
Producto final & Wireframe a mano + archivo index.html con estructura HTML semántica\\
\end{tabularx}
\normalsize

\puente{Antes de empezar, mira el viaje completo. Hoy le construyes \emph{la casa} a tu marca --- una página web con estructura semántica clara, lista para crecer en las guías siguientes.}

\begin{titlebox}{milcGradeDark}
Ruta MILC: así vamos a aprender
\end{titlebox}
\begin{tabularx}{\textwidth}{>{\centering\arraybackslash}X>{\centering\arraybackslash}X>{\centering\arraybackslash}X>{\centering\arraybackslash}X}
\phasecard{1}{milcVerde}{milcGris}{Escuta\\[-1mm]\small Observo, escucho y pregunto.} &
\phasecard{2}{milcTurquesa}{milcGris}{Sistematización\\[-1mm]\small Pensamiento computacional.} &
\phasecard{3}{milcMagenta}{milcGris}{Praxis\\[-1mm]\small Construyo y aplico.} &
\phasecard{4}{milcVino}{milcGris}{Evaluación\\[-1mm]\small Reflexión liberadora.}
\end{tabularx}


\puente{Antes de construir tu propia casa digital, vas a aprender a leer las que ya existen. Detrás de cada sitio bien hecho hay una estructura semántica --- visible solo si sabes mirar.}

\begin{titlebox}{milcVerde}
Fase 1 · Escuta
\end{titlebox}

\begin{softbox}{milcVerde}{milcGris}{Escena para escuchar}
\textbf{Actividad 1 · IDENTIFICA --- Lee la estructura semántica de 3 sitios} (15 min · individual). Visita 3 sitios distintos: un periódico, una página de universidad, un sitio de empresa. En cada uno: ¿qué hay en el header (logo, menú)? ¿qué en el main (contenido principal)? ¿qué en el footer (créditos, contacto)? Toma capturas y márcalas a lápiz si puedes.
\end{softbox}

\checkbox Visité un sitio de noticias y reconocí header, main y footer.\\
\checkbox Visité un sitio de universidad y reconocí su nav principal.\\
\checkbox Visité un sitio de empresa y reconocí su jerarquía visual.

\begin{infoband}{milcVerde}
\textbf{Lo que va al cuaderno (Actividad 1 · IDENTIFICA).} Título: «Actividad 1 --- Estructura semántica de 3 sitios». \textbf{Formato:} tabla de 4 columnas (Sitio | Header | Main | Footer), con qué contiene cada zona en una frase corta. \textbf{Extensión:} 3 filas, una por sitio. \textbf{Sabes que terminaste cuando} cada celda nombra elementos concretos (no ``varias cosas'') y reconoces el patrón común a los 3.
\end{infoband}

\puente{Ya viste cómo se organizan los sitios reales. Ahora vamos a aprender el \textbf{vocabulario exacto} que usa HTML para esa organización. Cinco o seis etiquetas semánticas resuelven el 90\% de los casos.}

\begin{titlebox}{milcTurquesa}
Fase 2 · Sistematización con pensamiento computacional
\end{titlebox}

\begin{softbox}{milcTurquesa}{milcGris}{Cuatro pilares aplicados al tema}
HTML semántico organiza una página en bloques con significado. Vamos a descomponerlo con los pilares del pensamiento computacional para que sepas qué etiqueta usar en cada lugar.
\end{softbox}

\small
\begin{tabularx}{\textwidth}{>{\bfseries\RaggedRight\arraybackslash}p{3.6cm}Y}
\rowcolor{milcTurquesa!90}\color{white}Pilar & \color{white}\bfseries Aplicación al tema\\
1. Descomposición & Una página se descompone en bloques semánticos: \texttt{header} (cabecera con logo y navegación), \texttt{nav} (menú), \texttt{main} (contenido único de la página), \texttt{section} (secciones temáticas), \texttt{article} (contenido autónomo), \texttt{aside} (información complementaria), \texttt{footer} (pie).\\
2. Reconocimiento de patrones & Todo sitio sigue un patrón: header arriba, footer abajo, main en el centro con secciones. Si el patrón se rompe sin razón, el usuario se pierde. La buena variación es excepción, no regla.\\
3. Abstracción & Lo esencial de una página cabe en su esqueleto semántico: si tu HTML sin estilos todavía se lee como un documento ordenado (h1, h2, párrafos, listas), tu abstracción está bien.\\
4. Algoritmo conceptual & Pasos: (1) dibuja a mano la jerarquía (header/nav/main/footer); (2) decide qué \texttt{section}s irán dentro del main; (3) traduce a HTML semántico; (4) prueba el resultado SIN CSS --- si se entiende, hiciste bien.\\
\end{tabularx}
\normalsize

\begin{drawbox}{milcTurquesa}{Anatomía de una página HTML semántica}
\texttt{header}: logo + nav (cabecera del sitio). \\[1mm] \texttt{nav}: menú principal (puede ir dentro de header). \\[1mm] \texttt{main}: contenido único de esta página. \\[1mm] \texttt{section}: agrupaciones temáticas dentro de main. \\[1mm] \texttt{article}: contenido autónomo (post, ficha, item). \\[1mm] \texttt{aside}: complemento lateral o relacionado. \\[1mm] \texttt{footer}: créditos, contacto, navegación secundaria.
\end{drawbox}

\begin{softbox}{milcMostaza}{milcGris}{Errores comunes y su corrección}
(1) Usar \texttt{div} para todo: pierdes semántica. (2) Múltiples \texttt{h1} por página: confunde a lectores de pantalla. (3) \texttt{main} repetido: solo debe haber UNO por página. (4) Confundir \texttt{section} con \texttt{article}: section agrupa, article es autónomo.
\end{softbox}

\begin{infoband}{milcTurquesa}
\textbf{Lo que va al cuaderno (Actividad 2 · EXPLICA).} Título: «Actividad 2 --- Vocabulario semántico HTML». \textbf{Formato:} 7 fichas, una por etiqueta (header / nav / main / section / article / aside / footer), cada ficha con nombre + para qué sirve (1 frase con tus palabras) + un ejemplo concreto de tu propio sitio. \textbf{Sabes que terminaste cuando} puedes explicar cada etiqueta sin mirar notas y dar un ejemplo específico de cómo la usarás.
\end{infoband}

\puente{Tienes el vocabulario y los referentes. Toca decidir: ¿qué secciones tendrá tu sitio y cómo se distribuirán? Vas a dibujar primero a mano (wireframe), después convertir a HTML.}

\begin{titlebox}{milcMagenta}
Fase 3 · Praxis con pensamiento computacional
\end{titlebox}

\begin{softbox}{milcMagenta}{milcGris}{Construyendo el producto con los cuatro pilares}
\textbf{Actividad 3 · CREA --- Wireframe + HTML semántico de tu home} (35 min · individual). Hora de construir. Primero dibujas a mano (wireframe), después conviertes a HTML real. El resultado es tu \texttt{index.html} con la estructura semántica de la página de inicio de tu sitio personal.
\end{softbox}

\small
\begin{tabularx}{\textwidth}{>{\bfseries\RaggedRight\arraybackslash}p{3.6cm}Y}
\rowcolor{milcMagenta!90}\color{white}Pilar & \color{white}\bfseries Aplicación al producto\\
Descomposición & Tu home se descompone en 4 zonas: header (nombre + nav), main (presentación + 2-3 secciones), aside opcional, footer (contacto + redes).\\
Patrones & Sigue el patrón: header arriba con tu nombre y un nav corto (3-5 enlaces), main con un h1 claro y secciones bien delimitadas, footer al final con email y redes profesionales.\\
Abstracción & Cada sección expresa una sola idea. Si tu \texttt{section} contiene 3 temas distintos, pártela. Si una sección es muy corta, fusiónala.\\
Algoritmo de armado & Algoritmo: (1) dibuja a mano la estructura en una hoja; (2) abre un archivo \texttt{index.html} en VS Code; (3) escribe el esqueleto semántico vacío; (4) llena cada bloque con texto real (sin CSS, todavía); (5) abre el HTML en Chrome y revisa.\\
\end{tabularx}
\normalsize

\begin{drawbox}{milcMagenta}{Checklist del HTML antes de cerrar el cuaderno}
\checkbox Wireframe a mano con header / main / footer claramente delimitados. \\[1mm] \checkbox Archivo \texttt{index.html} con doctype, html, head y body completos. \\[1mm] \checkbox Un solo \texttt{h1} por página. \\[1mm] \checkbox Etiquetas semánticas en lugar de \texttt{div} cuando aplica. \\[1mm] \checkbox Sin CSS, la página todavía se lee como un documento ordenado en Chrome.
\end{drawbox}

\begin{softbox}{milcMagenta}{milcGris}{Plantilla de trabajo}
\textbf{Header.} \textit{Tu nombre o lockup + nav corto.} \\[1mm] \textbf{Main / Hero.} \textit{h1 con tu propuesta de valor + breve presentación.} \\[1mm] \textbf{Main / Sobre mí.} \textit{section con 1 párrafo personal.} \\[1mm] \textbf{Main / Proyectos o servicios.} \textit{section con 2-3 articles.} \\[1mm] \textbf{Footer.} \textit{Email, redes profesionales, derechos.}
\end{softbox}

\begin{infoband}{milcMagenta}
\textbf{Lo que va al cuaderno (Actividad 3 · CREA).} Título: «Actividad 3 --- Mi home semántico v1». \textbf{Formato:} (1) wireframe a mano en la página izquierda, (2) transcripción del HTML real en la página derecha. \textbf{Extensión:} dos páginas enfrentadas. \textbf{Sabes que terminaste cuando} tu HTML carga en Chrome sin errores, tiene un solo h1, y sin CSS la página se entiende como documento estructurado.
\end{infoband}

\puente{Lo que sigue resume \emph{qué} entregas hoy: un wireframe a mano y un archivo HTML. El criterio principal: que un extraño pueda leer tu HTML sin estilos y entender la página.}

\begin{titlebox}{milcMostaza}
Producto de la sesión
\end{titlebox}
\textbf{Wireframe a mano + archivo index.html con estructura HTML semántica}.
\begin{softbox}{milcMostaza}{milcGris}{Criterios de calidad}
(1) El wireframe muestra las 3 zonas (header / main / footer) claramente. (2) El HTML usa al menos 5 etiquetas semánticas distintas. (3) Existe un único \texttt{h1}. (4) Sin CSS, el navegador todavía muestra una jerarquía clara. (5) El HTML pasa el validador W3C sin errores estructurales.
\end{softbox}

\puente{Antes de cerrar, mira tu sitio desde las cinco dimensiones humanas. Una página web también incluye o excluye personas; piensa en quién \emph{no} podría leerla y por qué.}

\begin{titlebox}{milcVino}
Fase 4 · Evaluación liberadora — cinco dimensiones
\end{titlebox}

\begin{softbox}{milcVino}{milcGris}{Sentido de la evaluación}
Evaluar no es solo revisar una nota. Es reconocer qué aprendí, qué cambió en mí, qué emociones manejé, cómo me relaciono con la ciudad y la comunidad, y qué le dejaré a quien viene detrás.
\end{softbox}

\textbf{Mi reflexión liberadora en cinco dimensiones.} Completa cada frase iniciada:

\small
\begin{tabularx}{\textwidth}{>{\bfseries\RaggedRight\arraybackslash}p{3.4cm}Y>{\RaggedRight\arraybackslash}p{4.6cm}}
\rowcolor{milcVino}\color{white}Dimensión & \color{white}\bfseries Mi respuesta (frame starter) & \color{white}\bfseries Algo que descubrí\\
1. Personal & Lo que más cambió en mí fue \linead{6.5cm} & \linead{4.2cm}\\
2. Emocional & La emoción más fuerte fue \linead{2.5cm} y la regulé \linead{3cm} & \linead{4.2cm}\\
3. Ciudadana & En democracia esto importa porque \linead{6cm} & \linead{4.2cm}\\
4. Local (Cartago) & En mi barrio sería útil para \linead{6.5cm} & \linead{4.2cm}\\
5. Intergeneracional & A mi abuela/abuelo le diría \linead{6.5cm} & \linead{4.2cm}\\
\end{tabularx}
\normalsize

\begin{infoband}{milcVino}
\textbf{Para la próxima sustentación.} Vuelve a esta tabla en seis meses. Verás que las respuestas cambian — no porque lo aprendido cambie, sino porque tú cambias. La evaluación liberadora es una conversación contigo a través del tiempo.
\end{infoband}


\puente{Y para terminar, tres voces sobre las casas digitales. Dussel sobre quién entra y quién no, Séneca sobre el orden como disciplina, Floridi sobre la web como espacio común.}

\begin{titlebox}{milcGradeDark}
Triángulo de pensamiento · cierre filosófico
\end{titlebox}

\begin{softbox}{milcVino}{milcGris}{Dussel · lente del nosotros}
\textit{``La analéctica supera a la dialéctica al partir del Otro como Otro.''}

Tu sitio web es una puerta. Cada decisión de diseño define a quién deja entrar y a quién deja afuera. La accesibilidad (etiquetas semánticas, contraste, texto alternativo) no es un ``extra técnico'': es la diferencia entre invitar al Otro o ignorarlo.

\textbf{¿Quién no puede leer mi sitio? ¿Una persona con lector de pantalla, con conexión lenta, que no maneja bien el español, lo logra navegar?}
\end{softbox}
\linead{\linewidth}\\[1mm]
\linead{\linewidth}

\begin{softbox}{milcMostaza}{milcGris}{Estoicismo · Séneca · cuidado interior}
\textit{``Nada se construye bien sin orden; sin él, la abundancia es ruina.''}

Un sitio web mal estructurado se siente ``hecho a la carrera''. Ordenar el HTML es un acto de disciplina --- pocas decisiones, en el orden correcto --- y depende solo de ti. Nadie te ordena el código; tú decides.

\textbf{¿Mi sitio refleja una estructura ordenada y deliberada, o se siente improvisado? ¿Qué dice eso de mi disciplina?}
\end{softbox}
\linead{\linewidth}\\[1mm]
\linead{\linewidth}

\begin{softbox}{milcTurquesa}{milcGris}{Floridi · lente de la infoesfera}
\textit{``La web es un patrimonio informacional común; cada sitio bien construido lo enriquece, cada sitio mal hecho lo erosiona.''}

Cuando publicas tu sitio, te sumas a la infoesfera pública. Tu HTML será leído por buscadores, lectores de pantalla, archivos del futuro. Construirlo con cuidado es ética informacional aplicada --- el equivalente moderno a sembrar bien.

\textbf{¿Mi sitio enriquece o ensucia el patrimonio digital común? ¿Estoy publicando algo de lo que el internet del futuro se beneficiará?}
\end{softbox}
\linead{\linewidth}\\[1mm]
\linead{\linewidth}

\begin{titlebox}{milcVino}
Compromiso final
\end{titlebox}

\small
\begin{tabularx}{\textwidth}{>{\RaggedRight\arraybackslash}p{3.0cm}YYY}
\rowcolor{milcVino}\color{white}\bfseries Criterio & \color{white}\bfseries Lo logré & \color{white}\bfseries En proceso & \color{white}\bfseries Necesito apoyo\\
Comprensión del tema & Explico ideas centrales con mis palabras. & Reconozco ideas pero falta precisión. & Necesito repasar conceptos base.\\
Pensamiento computacional & Apliqué los 4 pilares al diseñar mi producto. & Apliqué algunos pilares. & Necesito releer la fase 2.\\
Aplicación y producto & Mi producto cumple los criterios y se puede revisar. & Producto funcional con ajustes pendientes. & Aún en construcción.\\
Reflexión 5 dimensiones & Respondí cada dimensión con honestidad. & Respondí algunas con detalle. & Mis respuestas son muy generales.\\
Triángulo de pensamiento & Conecté las 3 voces con mi trabajo real. & Conecté al menos una. & No relaciono las voces con mi proceso.\\
\end{tabularx}
\normalsize

\textbf{Hoy entendí que:}\\[1mm]
\linead{\linewidth}\\[1mm]
\linead{\linewidth}

\textbf{Mi compromiso para la próxima sesión es:}\\[1mm]
\linead{\linewidth}\\[1mm]
\linead{\linewidth}

\begin{softbox}{milcMostaza}{milcGris}{Cita de despedida · Marco Aurelio}
\textit{``El obstáculo en el camino se vuelve el camino. Lo que estorba al camino se vuelve camino.''}\\[1mm]
Cada error de hoy, bien observado, se convierte en pieza del camino que pisarás mañana.
\end{softbox}

\small
\begin{tabularx}{\textwidth}{>{\bfseries}p{3.6cm}Y}
\rowcolor{milcGradeDark!12}Nombre completo: & \linead{10cm}\\
Fecha: & \linead{4cm} \quad Curso/grupo: \linead{4cm}\\
Mi firma: & \linead{9cm}\\
\end{tabularx}
\normalsize

\begin{infoband}{milcGradeDark}
\textbf{Para la próxima semana.} Comparte tu \texttt{index.html} con un compañero por correo. Pídele que lo abra en Chrome y te diga, en 1 frase, qué entendió. Si esa frase coincide con lo que querías comunicar, tu estructura semántica funciona. El próximo lunes traes ese feedback al aula.
\end{infoband}

\end{document}
